\section{External Validation: CADE Cartel Convictions}
\label{sec:cade}

We validate the FL screen against CADE (Conselho Administrativo de
Defesa Econ\^omica) cartel conviction records (2009--2019). Of 65
CADE procurement cartel cases, 47 convicted firms match to BEC via
CNPJ. Among 2,735 FL firms, 193 (7.1\%) co-participate with at least
one CADE-convicted firm in the same tender-item. A permutation test
(1,000 iterations, stratified by participation quartile) yields a
permutation mean of 2.0\% (SD~$= 0.3$\%), implying that the observed
7.1\% is 3.5 times higher than expected ($p < 0.001$). That 92.9\%
of FL firms have no CADE overlap is consistent with the low base rate
of cartel prosecution: if CADE captures only 1--5\% of actual cartels
\citep{connor2007price}, the 7.1\% overlap rate implies a much larger
fraction involved in undetected collusion.

\paragraph{Interpretation of co-participation.}
Co-participation---sharing a tender-item with a convicted firm---measures
\emph{market proximity} to known cartels, not direct cartel membership.
An FL firm that co-participates with a convicted cartel winner may be a
cover bidder deployed by that cartel, an unwitting competitor in
cartel-affected markets, or an unrelated participant that happens to
overlap. The permutation test controls for participation volume, and
the 3.5$\times$ excess rate indicates that FL firms are
disproportionately present in cartel-affected markets---consistent with
the screening hypothesis (H1). It does not prove that every
co-participating FL firm is a cartel member.

\paragraph{Dual pattern.}
The matching reveals a dual pattern consistent with the structural
model. The majority of CADE-convicted firms in BEC are
\emph{winners}---the designated recipients of cartel rents. However,
three convicted firms are themselves classified as frequent losers:
(i)~Sol Tecnologia em Energias Renov\'aveis (solar water heater
cartel, CADE process 08012.001273/2010-24; 84 BEC participations,
zero wins); (ii)~Nova Esperan\c{c}a Locadora de Ve\'iculos (school
transportation cartel, process 08700.005876/2019-85; 97
participations, zero wins); and (iii)~Jofran Com\'ercio de Produtos
para Higieniza\c{c}\~ao (trash bag cartel, Opera\c{c}\~ao Colludium,
process 08700.005789/2015-02; 65 participations, zero wins). All three
participated across multiple item groups and years without ever
winning---precisely the behavior that defines frequent losers.

\paragraph{Co-bidder analysis.}
FL firms appear in 5.6\% of CADE-firm tenders, compared to a baseline
of 2.1\% across all BEC tenders---a ratio of 2.6$\times$. For specific
convicted firms, the FL co-bidding rate is substantially higher:
69.3\% for Mayfran (school transportation), 68.8\% for Bosch
Termotecnologia (solar heaters), 55.2\% for Tejofran (trains/metro),
and 31.7\% for Convida (school meals). This co-presence pattern
validates the screen through two complementary channels: direct
identification of some cartel members and detection of the
operational infrastructure of bid-rigging schemes.

\paragraph{Robustness: excluding directly convicted FL firms.}
Three of the convicted firms are themselves classified as FL (Sol
Tecnologia, Nova Esperan\c{c}a, Jofran; see Dual Pattern above).
Excluding these three firms and their 246 associated tender-items from
the baseline regression leaves $\hat{\beta}$ unchanged at 0.064,
confirming that the price association is not driven by the small
number of FL firms with direct CADE involvement.

\paragraph{Beyond known cartels.}
Dropping all 31,447 CADE-involved tender-items and re-estimating the
baseline price regression yields $\hat{\beta} = 0.062$ ($N =
1{,}622{,}954$), virtually identical to the full-sample 0.064
(Table~\ref{tab:excl_cade}), demonstrating that FL detects price
anomalies \emph{beyond} known cartels.

